% =====================================================================
%  CAPÍTULO 3 — FUNDAMENTAÇÃO TEÓRICA
% =====================================================================

\chapter{FUNDAMENTAÇÃO TEÓRICA}
\label{cap:fundamentacao}

Este capítulo apresenta os conceitos que sustentam o \argus{}. Ele não é um
catálogo de tecnologias: procura explicar por que cada elemento existe, qual
problema resolve e onde termina a sua validade. A exposição está organizada em
seis blocos: comunicação com equipamentos, formatos de percurso e de atividade,
modelo físico da simulação, fisiologia e carga de treino, engenharia de
software e inteligência artificial local. A distinção entre \emph{medir},
\emph{estimar} e \emph{inferir} atravessa o capítulo e retorna na síntese final,
porque ela é central para a honestidade das afirmações do trabalho.

\section{Comunicação e controle de equipamentos}
\label{sec:ble-gatt}

\subsection{Bluetooth Low Energy e o modelo GATT}
\label{sec:ble-gatt-modelo}

O \textit{Bluetooth Low Energy} (BLE) é uma tecnologia de comunicação sem fio
de curto alcance voltada a dispositivos que transmitem pequenas quantidades de
dados de forma intermitente, com ênfase em baixo consumo de energia
\cite{bluetoothcore2023}. No ciclismo virtual, sensores de frequência cardíaca,
medidores de potência, sensores de cadência e rolos inteligentes usam BLE para
entregar telemetria e, no caso dos rolos, para receber comandos de resistência.

O \textit{Generic Attribute Profile} (GATT) organiza os dados em serviços e
características. Um serviço é identificado por um UUID e agrupa características
relacionadas; cada característica tem um valor e pode ser lida, escrita ou
notificada. Quando uma característica permite notificações, o cliente habilita
o descritor de configuração (CCCD) e passa a receber atualizações sempre que o
periférico publica um novo valor \cite{bluetoothcore2023}. O modelo é
assimétrico: o periférico mantém o estado e o cliente BLE, no caso o
computador ou o telefone, assina as notificações.

Três serviços são diretamente relevantes para o \argus{}:

\begin{itemize}
    \item \textbf{Heart Rate Service} (\texttt{0x180D}): expõe a frequência
        cardíaca instantânea, em batimentos por minuto. A característica
        \texttt{0x2A37} admite valores de 8 ou 16 bits e traz bits de estado
        que indicam o formato e a presença de informações adicionais.
    \item \textbf{Cycling Power Service} (\texttt{0x1818}): transmite potência
        instantânea em watts e, quando disponível, cadência calculada a partir
        das revoluções do pedivela. A característica \texttt{0x2A63} usa um
        campo de \textit{flags} para indicar quais grandezas opcionais estão
        presentes, o que exige leitura posicional e respeito à ordem dos
        campos.
    \item \textbf{Fitness Machine Service} (FTMS, \texttt{0x1826}): padroniza
        a leitura de dados de equipamentos de fitness e o envio de comandos de
        controle, como potência alvo e inclinação simulada
        \cite{bluetooth2017}.
\end{itemize}

O BLE não garante uma taxa de amostragem fixa. O intervalo entre notificações
depende do dispositivo, do firmware, do intervalo de conexão e das condições de
rádio. Essa característica é importante para o \argus{}: o sistema registra cada
notificação recebida e usa o intervalo entre chegadas para integrar a distância
percorrida; entretanto, algumas métricas agregadas contam amostras como se cada
uma representasse um segundo. A suposição é razoável para os equipamentos
testados, que enviam dados aproximadamente a \SI{1}{\hertz}, mas não é uma
garantia do protocolo. O Capítulo~\ref{cap:implementacao} mostra como essa
diferença aparece no código e o Capítulo~\ref{cap:resultados} a registra como
limitação.

\subsection{FTMS: um caminho padronizado de controle}
\label{sec:ftms}

O FTMS define características para recursos do equipamento, faixa de potência
suportada, dados de bicicleta do perfil \textit{Indoor Bike Data}, status e ponto
de controle. O ponto de controle (\texttt{0x2AD9}) recebe comandos e responde
por indicação; o cliente deve solicitar o controle antes de enviar comandos de
operação. Entre os
opcodes relevantes estão a requisição de controle (\texttt{0x00}), o ajuste de
resistência (\texttt{0x04}), a potência alvo (\texttt{0x05}) e a inclinação
alvo (\texttt{0x03}) \cite{bluetooth2017}.

O \argus{} usa o FTMS em dois modos. No modo \emph{ERG}, envia a potência alvo;
no modo \emph{simulação}, envia a inclinação calculada a partir do GPX. O
equipamento interpreta o comando e ajusta internamente o freio; o software não
controla diretamente a corrente elétrica do rolo. Essa separação é importante:
o modelo físico do \argus{} calcula a velocidade virtual, mas a resistência
efetivamente aplicada depende do firmware do rolo, da calibração e da
resposta dinâmica do equipamento. O FTMS padroniza a mensagem, não o
comportamento mecânico de cada modelo.

\subsection{FE-C sobre BLE: compatibilidade sem ANT+ nativo}
\label{sec:fec-ble}

O \textit{Fitness Equipment Control} (FE-C) é um perfil da rede ANT+ para
controle de equipamentos de fitness. Ele define páginas de dados para potência
alvo, resistência de pista, resistência do vento, configuração do usuário e
telemetria do rolo \cite{antfec2024}. O ANT+ nativo utiliza o rádio ANT e,
em computadores, costuma exigir um adaptador USB específico. O \argus{} não
implementa esse caminho nativo.

O que o \argus{} implementa é o transporte das mensagens FE-C sobre BLE, por
meio de um serviço vendor-specific exposto por alguns rolos. O serviço
usa o UUID \texttt{6e40fec1-b5a3-f393-e0a9-e50e24dcca9e}; a característica de
notificação é \texttt{6e40fec2-...} e a de escrita é \texttt{6e40fec3-...}.
Cada mensagem tem 13 bytes: sincronismo (\texttt{0xA4}), comprimento
(\texttt{0x09}), identificador (\texttt{0x4F}), canal (\texttt{0x05}), página,
sete bytes de carga útil e um \textit{checksum} XOR. As páginas usadas pelo
sistema estão resumidas na Tabela~\ref{tab:fec-paginas}.

\begin{table}[htbp]
    \centering
    \caption{Páginas FE-C usadas pelo Argus}
    \label{tab:fec-paginas}
    \small
    \begin{tabular}{@{}>{\raggedright\arraybackslash}p{1.4cm}
        >{\raggedright\arraybackslash}p{3.2cm}
        >{\raggedright\arraybackslash}p{8.6cm}@{}}
        \toprule
        \textbf{Página} & \textbf{Função} & \textbf{Uso e escala no Argus} \\
        \midrule
        49 & Potência alvo & Modo ERG; resolução de \SI{0.25}{\watt} \\
        51 & Resistência de pista & Modo simulação; inclinação com resolução
        de \SI{0.01}{\percent} e deslocamento de \SI{200}{\percent};
        coeficiente de rolamento com resolução de 0,00005 \\
        55 & Configuração do usuário & Peso do ciclista e da bicicleta;
        enviado na inicialização com valores fixos de \SI{75}{\kilogram} e
        \SI{10}{\kilogram} \\
        50 & Resistência do vento & Inicialização com resistência de 0,51 e
        fator de drafting 1,0 \\
        25 & Dados específicos do rolo & Leitura de potência (12 bits) e
        cadência nas notificações \\
        \bottomrule
    \end{tabular}
    \vspace{0.2cm}
    \par\small Fonte: autoria própria, com base na especificação FE-C
    \cite{antfec2024} e no código do projeto.
\end{table}

A Tabela~\ref{tab:fec-paginas} também revela um limite de implementação: a
página 55 é enviada com valores fixos, e não com o peso do perfil ativo. O peso
do ciclista entra no motor físico do \argus{}, mas não é reenviado ao rolo
quando o perfil muda. Em rolos de treinamento que usam a configuração de
usuário para ajustar a resistência, essa diferença pode afetar a sensação de
esforço. O trabalho registra o fato; a correção não faz parte do escopo desta
versão.

\begin{table}[htbp]
    \centering
    \caption{Comparação entre os dois caminhos de controle usados no desktop}
    \label{tab:ftms-fec}
    \small
    \begin{tabular}{@{}>{\raggedright\arraybackslash}p{2.6cm}
        >{\raggedright\arraybackslash}p{4.8cm}
        >{\raggedright\arraybackslash}p{5.8cm}@{}}
        \toprule
        \textbf{Aspecto} & \textbf{FTMS} & \textbf{FE-C sobre BLE} \\
        \midrule
        Origem & Padrão Bluetooth SIG & Perfil ANT+ transportado em serviço
        vendor-specific \\
        UUID & \texttt{0x1826} & \texttt{6e40fec1-...} \\
        Controle & Ponto de controle \texttt{0x2AD9} & Páginas 49 e 51 \\
        Telemetria & \textit{Indoor Bike Data} & Página 25 e/ou Cycling Power \\
        No Argus & Desktop & Desktop e Android \\
        Limite & Variação entre fabricantes & Não é ANT+ nativo; depende do
        serviço proprietário \\
        \bottomrule
    \end{tabular}
    \vspace{0.2cm}
    \par\small Fonte: autoria própria.
\end{table}

A coexistência dos dois caminhos amplia a compatibilidade, mas não deve ser
descrita como suporte a ANT+. O \argus{} fala FE-C \emph{sobre BLE}; um
adaptador ANT+ USB, um Garmin Edge ou um rolo que exponha apenas ANT+ nativo
não são reconhecidos por esse código. Essa distinção é relevante para a
reprodutibilidade: o leitor precisa saber qual transporte foi testado.

\subsection{Leitura, escalas e taxas de amostragem}
\label{sec:leitura-ble}

A leitura de características BLE exige interpretar \textit{flags}, respeitar a
ordem dos campos e aplicar as escalas do perfil. No Cycling Power Service, por
exemplo, a potência ocupa dois bytes em \textit{little-endian}; a cadência só
está disponível quando o bit correspondente está ligado e é calculada a partir
da diferença de revoluções e de tempo entre duas notificações. No FE-C, a
potência é um campo de 12 bits e a cadência é lida diretamente de um byte. Um
erro de deslocamento não gera uma mensagem inválida: gera um número plausível e
errado, o que é mais difícil de detectar.

O \argus{} adota três precauções: valida o tamanho mínimo do pacote, verifica
os bits de presença antes de ler campos opcionais e usa o instante de chegada
da notificação como \textit{timestamp}. Ainda assim, não há reamostragem
explícita para uma grade uniforme de \SI{1}{\second}. O intervalo entre
notificações é usado no cálculo de distância, mas os vetores de potência e de
frequência cardíaca são formados por amostras consecutivas. Se um dispositivo
enviar duas notificações por segundo, as métricas que somam amostras como
segundos superestimarão o tempo; se enviar uma a cada dois segundos,
subestimarão. A limitação é tratada novamente nos Capítulos~\ref{cap:implementacao}
e~\ref{cap:resultados}.

\section{Formatos de percurso, atividade e treino}
\label{sec:formatos}

\subsection{GPX: geometria e altimetria}
\label{sec:gpx}

O \textit{GPS Exchange Format} (GPX) é um formato XML aberto para intercâmbio
de pontos geográficos, rotas e trilhas \cite{gpx2004}. Um arquivo pode conter
\textit{waypoints}, rotas ou \textit{tracks}; cada ponto tem latitude, longitude
e, opcionalmente, elevação e tempo. O \argus{} aceita \textit{tracks} e, quando
não há trilha, tenta ler rotas. A elevação é a informação que alimenta o perfil
altimétrico e o cálculo de inclinação.

Há duas decisões geométricas importantes. A primeira é o cálculo da distância:
o desktop usa a distância tridimensional entre pontos consecutivos, incluindo
a diferença de elevação; o Android usa a distância horizontal calculada pela
fórmula de Haversine, sem elevação. A segunda é a suavização: o desktop calcula
a inclinação em uma janela de até cinco pontos para cada lado e limita o
resultado a $\pm 25\%$; o Android usa uma janela de dois pontos para cada lado e
não aplica o limite. A Tabela~\ref{tab:gpx-diferencas} resume a diferença.

Essas diferenças não são detalhes irrelevantes. A suavização maior no desktop
reduz picos espúrios de elevação, mas também atenua rampas curtas; a janela
menor no Android preserva mais detalhe e mais ruído. A interpolação do desktop
produz uma posição contínua ao longo do segmento; o Android retorna o ponto
seguinte, o que torna a inclinação e a posição discretas. Para o uso de
reconhecimento, a diferença é aceitável; para comparar sessões entre
plataformas, ela precisa ser considerada.

\enlargethispage{\baselineskip}
O GPX tem limites que nenhum processamento elimina: a elevação pode vir de
GPS, barômetro ou modelo digital de terreno, com ruído e descontinuidades; o
arquivo não descreve superfície, clima, trânsito, curvas ou regras de trânsito;
e a densidade de pontos varia. O \argus{} usa o GPX como representação
aproximada do relevo, não como réplica da estrada.

\begin{table}[H]
    \centering
    \caption{Pipeline de GPX no desktop e no Android}
    \label{tab:gpx-diferencas}
    \small
    \begin{tabular}{@{}>{\raggedright\arraybackslash}p{3.4cm}
        >{\raggedright\arraybackslash}p{5.2cm}
        >{\raggedright\arraybackslash}p{5.2cm}@{}}
        \toprule
        \textbf{Etapa} & \textbf{Desktop (Go)} & \textbf{Android (JavaScript)} \\
        \midrule
        Distância & Tridimensional, com elevação & Horizontal (Haversine), sem
        elevação \\
        Inclinação & Janela de $\pm 5$ pontos & Janela de $\pm 2$ pontos \\
        Limite & $\pm 25\%$ & Sem limite explícito \\
        Consulta por distância & Interpolação linear entre pontos & Primeiro
        ponto com distância maior ou igual \\
        Percurso & Repetição modular ao atingir o fim & Repetição modular ao
        atingir o fim \\
        \bottomrule
    \end{tabular}
    \vspace{0.2cm}
    \par\small Fonte: autoria própria, com base em
    \code{internal/service/gpx/service.go} e \code{frontend/src/main.js}.
\end{table}

\subsection{FIT: registro binário da atividade}
\label{sec:fit}

O \textit{Flexible and Interoperable Data Transfer} (FIT) é um formato binário
mantido pela Garmin para registrar atividades e dados de sensores
\cite{garmin2024}. Um arquivo contém uma cabeçalho, mensagens de definição que
descrevem a estrutura dos registros e mensagens de dados; o cabeçalho e o
arquivo incluem verificações de integridade. O formato é compacto e
interoperável, mas não é autodescritivo como o XML: para interpretar um campo,
é preciso conhecer o perfil FIT e a mensagem correspondente.

Os valores usados pelo \argus{} seguem as escalas do perfil: latitude e
longitude em semicírculos, velocidade em milímetros por segundo, distância em
centímetros e altitude com escala 5 e deslocamento de 500 metros. A Tabela~
\ref{tab:fit-escalas} registra as conversões implementadas. O arquivo gerado
contém uma mensagem de identificação, os registros de telemetria, um evento de
parada, uma volta e uma sessão com totais. O subtipo é atividade virtual, o que
informa à plataforma de destino que o percurso não foi realizado em estrada.

\begin{table}[htbp]
    \centering
    \caption{Conversões de unidades na exportação FIT}
    \label{tab:fit-escalas}
    \small
    \begin{tabular}{@{}>{\raggedright\arraybackslash}p{4.2cm}
        >{\raggedright\arraybackslash}p{5.2cm}
        >{\raggedright\arraybackslash}p{4.0cm}@{}}
        \toprule
        \textbf{Grandeza} & \textbf{Unidade no Argus} & \textbf{Escala FIT} \\
        \midrule
        Latitude e longitude & graus decimais & semicírculos \\
        Velocidade & \si{\kilo\meter\per\hour} & \si{\milli\meter\per\second} \\
        Distância & metros & centímetros \\
        Altitude & metros & $(m + 500) \times 5$ \\
        Potência & watts & watts \\
        Cadência e frequência cardíaca & rpm e bpm & rpm e bpm \\
        \bottomrule
    \end{tabular}
    \vspace{0.2cm}
    \par\small Fonte: autoria própria, com base no perfil FIT
    \cite{garmin2024} e em \code{internal/service/fit/service.go}.
\end{table}

A aceitação do arquivo por Strava ou Intervals.icu é evidência de
interoperabilidade, não de exatidão fisiológica. O FIT transporta os valores
que o \argus{} calculou; se a potência, a distância ou a elevação de origem
estiverem erradas, o erro será preservado no arquivo. A conferência foi feita por
leitura nas plataformas de destino e por comparação manual de totais; a validação
formal com o validador oficial do FIT SDK permanece como trabalho futuro,
retomado no Capítulo~\ref{cap:resultados}.

\subsection{ZWO: planos de treino estruturados}
\label{sec:zwo}

O ZWO é o formato XML usado pelo Zwift para descrever treinos estruturados
\cite{zwift2024}. Cada passo tem duração e alvo de potência, geralmente expresso
como fração do FTP; há passos constantes, rampas, aquecimentos, resfriamentos e
intervalos repetidos. O formato é adequado ao modo ERG porque transforma o plano
em uma sequência de alvos de potência.

O \argus{} implementa um subconjunto do formato: lê passos
\textit{SteadyState}, \textit{Warmup}, \textit{Cooldown}, \textit{Ramp},
\textit{IntervalsT} e \textit{FreeRide}, e achata os intervalos repetidos em
segmentos individuais. Extensões não previstas são ignoradas. O assistente de
IA gera arquivos ZWO a partir de uma estrutura JSON e o mesmo carregador é usado
para planos externos e para planos gerados localmente. A vantagem é a
interoperabilidade com um formato conhecido; o limite é que a obediência
sintática ao ZWO não garante adequação fisiológica do treino. Um plano pode ser
válido e, ainda assim, inadequado para o atleta.

\section{Modelo físico da simulação}
\label{sec:equacao-martin}

\subsection{A equação de Martin}
\label{sec:equacao-martin-formal}

O movimento de uma bicicleta pode ser descrito pelo balanço entre a potência
aplicada pelo ciclista e as forças que se opõem ao avanço. O modelo de Martin
\cite{martin1998} reúne arrasto aerodinâmico, resistência ao rolamento,
componente gravitacional e aceleração:

\begin{equation}\label{eq:martin}
P_{\text{crank}} = \frac{1}{\eta}\left[
\frac{1}{2}\rho C_d A v^3
+ m g C_{rr}\cos\theta\,v
+ m g\sin\theta\,v
+ m a v
\right],
\end{equation}

\noindent em que $P_{\text{crank}}$ é a potência no pedivela, $\rho$ é a
densidade do ar, $C_dA$ é a área frontal efetiva multiplicada pelo coeficiente
de arrasto, $v$ é a velocidade, $m$ é a massa total, $g$ é a aceleração da
gravidade, $C_{rr}$ é o coeficiente de resistência ao rolamento, $\theta$ é o
ângulo de inclinação da via, $a$ é a aceleração linear e $\eta$ é a eficiência
da transmissão. O termo $\cos\theta$ aparece na resistência ao rolamento porque
a força normal diminui em rampas; o termo $\sin\theta$ representa a componente
do peso ao longo da via.

A implementação do \argus{} adota a hipótese quase estática: o termo de
aceleração é desprezado e a velocidade é obtida para cada potência e inclinação
instantâneas. A equação resolvida é, portanto,

\begin{equation}\label{eq:martin-estatica}
P_{\text{wheel}} = \frac{1}{2}\rho C_d A v^3
+ \left(m g C_{rr}\cos\theta + m g\sin\theta\right)v,
\end{equation}

\noindent com $P_{\text{wheel}} = \eta P_{\text{crank}}$. A Tabela~
\ref{tab:parametros-sim} lista os valores usados no código. Eles são
constantes: não variam com o perfil aerodinâmico, o tipo de pneu, a pressão, a
temperatura, a altitude ou a postura do ciclista.

\begin{table}[htbp]
    \centering
    \caption{Parâmetros fixos do motor de simulação}
    \label{tab:parametros-sim}
    \small
    \begin{tabular}{@{}>{\raggedright\arraybackslash}p{3.6cm}
        >{\raggedright\arraybackslash}p{3.4cm}
        >{\raggedright\arraybackslash}p{7.2cm}@{}}
        \toprule
        \textbf{Parâmetro} & \textbf{Valor} & \textbf{Observação} \\
        \midrule
        $g$ & \SI{9.81}{\meter\per\square\second} & Constante \\
        $\rho$ & \SI{1.225}{\kilogram\per\cubic\meter} & Ar ao nível do mar \\
        $C_dA$ & \SI{0.32}{\square\meter} & Valor fixo, não individualizado \\
        $C_{rr}$ & 0,005 & Desktop; o Android usa 0,004 \\
        $\eta$ & 0,96 & Desktop; o Android usa 0,97 \\
        Faixa de velocidade & 0 a \SI{40}{\meter\per\second} & Aproximadamente
        \SI{144}{\kilo\meter\per\hour} \\
        Iterações & 20 & Bisseção com tolerância de \SI{0.01}{\meter\per\second} \\
        \bottomrule
    \end{tabular}
    \vspace{0.2cm}
    \par\small Fonte: autoria própria, com base em
    \code{internal/service/sim/engine.go}.
\end{table}

\subsection{Solução numérica e estabilidade}
\label{sec:bissecao}

Para uma inclinação fixa, a função que relaciona potência e velocidade é
monótona crescente no intervalo fisicamente relevante: mais velocidade exige
mais potência. Essa monotonicidade permite usar bisseção. O algoritmo mantém um
intervalo $[0, 40]$ m/s, calcula a potência necessária no ponto médio e descarta
a metade que não contém a solução. Com 20 iterações, o intervalo inicial seria
reduzido a menos de $4 \times 10^{-5}$ m/s; na prática, o código interrompe
antes, quando a diferença entre os limites é inferior a \SI{0.01}{\meter\per\second}.
O método não divide pela velocidade, não depende de chute inicial e lida com
potência nula em descidas, situação em que a velocidade terminal é determinada
pelo equilíbrio entre gravidade, arrasto e rolamento.

A estabilidade numérica não equivale a realismo físico. O modelo não representa
a inércia do conjunto ciclista--bicicleta, a resposta dinâmica do rolo, a
folga da transmissão, a deformação do pneu ou a variação de $C_{rr}$ com o piso.
Também não há vento, \textit{drafting}, frenagem, curvas ou limite de aderência.
Em uma descida, o modelo pode produzir uma velocidade alta com potência zero,
mas não simula a necessidade de frear; em uma subida, a velocidade responde
instantaneamente à potência, sem o atraso que o conjunto mecânico e o volante do
rolo impõem. A simulação é, portanto, um modelo de primeira ordem para
explorar o relevo e comparar esforços, não um gêmeo digital da estrada.

\subsection{O que a velocidade significa}
\label{sec:velocidade-significado}

A velocidade exibida e a distância acumulada são saídas do modelo, não medidas
diretas de um sensor de velocidade. Elas dependem da potência informada pelo
medidor ou pelo rolo, da inclinação derivada do GPX e dos parâmetros fixos
da Tabela~\ref{tab:parametros-sim}. Consequências práticas:

\begin{itemize}
    \item a distância virtual pode divergir da distância registrada por um
        GPS em percurso real, sobretudo em trechos com vento, trânsito, curvas
        ou paradas;
    \item a comparação entre sessões é mais confiável quando o mesmo perfil,
        o mesmo rolo e o mesmo GPX são usados;
    \item a velocidade não deve ser usada como medida de desempenho absoluto
        em competições ou estudos que exijam validade externa;
    \item a potência, por outro lado, é medida pelo sensor (com a incerteza do
        próprio sensor) e é a grandeza mais adequada para comparar esforços.
\end{itemize}

Essa separação entre grandeza medida e grandeza modelada é uma das bases do
posicionamento do \argus{}: o sistema não promete reproduzir a estrada, promete
tornar explícito o modelo que usa para representá-la.

\section{Fisiologia e carga de treino}
\label{sec:fisiologia}

\subsection{Limiares, MAP e \texorpdfstring{VO\textsubscript{2}max}{VO2max}}
\label{sec:limiares}

A intensidade do exercício pode ser organizada em domínios e limiares. O
primeiro limiar (LT1 ou VT1) marca a transição em que o lactato começa a subir
acima do basal e a ventilação aumenta de forma desproporcional ao oxigênio; o
segundo (LT2, VT2 ou limiar anaeróbio) marca a intensidade mais alta em que o
lactato ainda se mantém em estado quase estacionário, próxima da máxima fase
estável de lactato (MLSS). Esses conceitos se relacionam, mas não são
sinônimos: LT1/LT2 são definidos por lactato; VT1/VT2, por variáveis
ventilatórias; MLSS, pela estabilidade do lactato \cite{faude2009,meyer2005}.

A potência aeróbia máxima (MAP) é a maior potência atingida em um teste
incremental, normalmente um teste de rampa. Ela é uma medida de potência, não
de consumo de oxigênio. O VO\textsubscript{2}max é a taxa máxima de captação de
oxigênio, medida por calorimetria indireta ou estimada por protocolos
específicos; não pode ser inferido apenas da potência de cinco minutos sem
um modelo de conversão validado para a população e o equipamento. O teste de
cinco minutos do \argus{} registra a melhor potência média de cinco minutos e
a usa como referência interna; ele é um teste de capacidade aeróbia usado como
proxy de MAP, não uma medida de VO\textsubscript{2}max. O rótulo da interface deve ser lido
com essa ressalva.



As linhas de VT1 e VT2 merecem uma explicação específica. Elas não são
calculadas a partir de trocas gasosas, lactato ou variabilidade da frequência
cardíaca. São marcadores visuais que o usuário pode arrastar sobre o gráfico
para indicar onde, na sua leitura, ocorreram as transições. A posição inicial
em 30\% e 70\% do eixo é apenas uma conveniência de interface. Portanto, o
software não ``detecta'' limiares; ele oferece uma ferramenta de anotação
manual. A diferença é grande e precisa ser declarada.

\begin{table}[htbp]
    \centering
    \caption{Conceitos de limiar e o tratamento dado pelo Argus}
    \label{tab:limiares}
    \small
    \begin{tabular}{@{}>{\raggedright\arraybackslash}p{2.8cm}
        >{\raggedright\arraybackslash}p{5.0cm}
        >{\raggedright\arraybackslash}p{6.4cm}@{}}
        \toprule
        \textbf{Conceito} & \textbf{O que representa} & \textbf{No Argus} \\
        \midrule
        VT1 / LT1 & Início do aumento sustentado de lactato e da ventilação &
        Duas linhas verticais móveis no gráfico, iniciadas em 30\% e 70\% do
        eixo temporal; não há detecção automática \\
        VT2 / LT2 & Transição para intensidade severa; próximo da MLSS &
        Mesmas linhas móveis; o usuário interpreta o gráfico \\
        MAP & Maior potência de um teste incremental & Teste de cinco minutos
        usado como proxy de MAP e de capacidade aeróbia \\
        VO\textsubscript{2}max & Captação máxima de oxigênio & Não medido e
        não estimado por equação própria \\
        FTP & Potência sustentável em estado quase estacionário por cerca de
        uma hora & Estimada por testes de rampa, 20 minutos ou 5 minutos, com
        fatores de conversão \\
        \bottomrule
    \end{tabular}
    \vspace{0.2cm}
    \par\small Fonte: autoria própria, com base em
    \cite{faude2009,meyer2005,coggan2010}.
\end{table}

\subsection{FTP e zonas de potência}
\label{sec:ftp-zonas}

A potência funcional limiar (FTP) é uma convenção de treino: a maior potência
que o atleta consegue sustentar em estado quase estacionário por aproximadamente
uma hora \cite{coggan2010}. Ela não é uma constante fisiológica e não é medida
diretamente em todos os testes. O \argus{} estima a FTP por três caminhos:
teste de rampa (melhor potência de um minuto multiplicada por 0,75), teste de
20 minutos (melhor potência de 20 minutos multiplicada por 0,95) e teste de
cinco minutos (melhor potência de cinco minutos multiplicada por 0,80). Os
fatores são convenções internas do aplicativo; não substituem um teste de FTP
bem conduzido nem uma medida de MLSS.

As zonas de potência implementadas são seis, inspiradas na escala de Coggan,
mas sem a sétima faixa:

\begin{table}[htbp]
    \centering
    \caption{Zonas de potência implementadas}
    \label{tab:zonas-potencia}
    \small
    \begin{tabular}{@{}>{\raggedright\arraybackslash}p{1.4cm}
        >{\raggedright\arraybackslash}p{3.8cm}
        >{\raggedright\arraybackslash}p{7.4cm}@{}}
        \toprule
        \textbf{Zona} & \textbf{Faixa em relação à FTP} & \textbf{Uso típico} \\
        \midrule
        Z1 & $< 55\%$ & Recuperação ativa \\
        Z2 & $55\%$ a $74\%$ & Resistência aeróbia \\
        Z3 & $75\%$ a $89\%$ & Tempo \\
        Z4 & $90\%$ a $104\%$ & Limiar \\
        Z5 & $105\%$ a $120\%$ & Capacidade aeróbia \\
        Z6 & $\geq 121\%$ & Potência anaeróbia e esforços curtos \\
        \bottomrule
    \end{tabular}
    \vspace{0.2cm}
    \par\small Fonte: autoria própria, com base em
    \code{internal/service/fit/metrics.go}.
\end{table}

A contagem de tempo nas zonas é feita percorrendo o vetor de potência e
incrementando a zona correspondente. O código pressupõe que cada amostra vale
um segundo; se a taxa de notificação for diferente, a contagem fica distorcida
na mesma proporção. Além disso, a zona é uma classificação arbitrária de uma
variável contínua: a fronteira entre Z3 e Z4 não corresponde a uma mudança
fisiológica abrupta.

\subsection{Zonas de frequência cardíaca}
\label{sec:zonas-fc}

As zonas de frequência cardíaca do \argus{} usam cinco faixas percentuais da
frequência cardíaca máxima cadastrada no perfil: abaixo de 60\%, 60--70\%,
70--80\%, 80--90\% e acima de 90\%. O campo LTHR (frequência cardíaca de
limiar de lactato) existe no perfil, mas não alimenta esse cálculo. Também não
são usados a reserva de frequência cardíaca (Karvonen) nem zonas derivadas de
teste de limiar.

\begin{table}[htbp]
    \centering
    \caption{Zonas de frequência cardíaca implementadas}
    \label{tab:zonas-fc}
    \small
    \begin{tabular}{@{}>{\raggedright\arraybackslash}p{1.6cm}
        >{\raggedright\arraybackslash}p{4.2cm}
        >{\raggedright\arraybackslash}p{6.8cm}@{}}
        \toprule
        \textbf{Zona} & \textbf{Faixa em relação à FC máxima} & \textbf{Rótulo
        na interface} \\
        \midrule
        Z1 & $< 60\%$ & Recuperação \\
        Z2 & $60\%$ a $69\%$ & Resistência \\
        Z3 & $70\%$ a $79\%$ & Tempo \\
        Z4 & $80\%$ a $89\%$ & Limiar \\
        Z5 & $\geq 90\%$ & VO\textsubscript{2}max \\
        \bottomrule
    \end{tabular}
    \vspace{0.2cm}
    \par\small Fonte: autoria própria, com base em
    \code{internal/service/fit/metrics.go}.
\end{table}

O rótulo ``VO\textsubscript{2}max'' na Z5 é uma convenção de treino, não uma
medida de consumo de oxigênio. A zona indica apenas que a frequência cardíaca
está acima de 90\% da máxima cadastrada. A precisão depende de a máxima estar
correta, do sensor e da relação individual entre frequência cardíaca e
intensidade. O método por percentual da máxima é simples e transparente, mas
menos individualizado do que zonas por LTHR, por reserva de frequência cardíaca
ou por limiares ventilatórios. A opção do \argus{} é defensável para
acompanhamento, não para prescrição clínica.

\subsection{Potência normalizada, fator de intensidade e TSS}
\label{sec:np-tss}

A potência normalizada (NP) procura representar o custo fisiológico de uma
sessão com intensidade variável. O cálculo aplica uma média móvel de 30
amostras, eleva cada média à quarta potência, calcula a média dessas potências
e extrai a raiz quarta:

\begin{equation}\label{eq:np}
NP = \left(\frac{1}{N}\sum_{i=1}^{N}\bar{p}_{30,i}^{4}\right)^{1/4},
\end{equation}

\noindent em que $\bar{p}_{30,i}$ é a i-ésima média móvel de 30 amostras e $N$
é o número de janelas. O expoente 4 faz com que picos de potência pesem mais
do que na média aritmética. A implementação usa amostras consecutivas; para
NP de 30 segundos, a amostragem precisa ser de aproximadamente \SI{1}{\hertz}.

O fator de intensidade (IF) e o \textit{Training Stress Score} (TSS) são
definidos por

\begin{equation}\label{eq:if}
IF = \frac{NP}{FTP},
\end{equation}

\begin{equation}\label{eq:tss}
TSS = \frac{d \cdot NP \cdot IF}{FTP \cdot 3600} \times 100,
\end{equation}

\noindent em que $d$ é a duração em segundos. Um TSS de 100 corresponde,
por definição, a uma hora no FTP. O TSS não é uma medida de calorias nem de
dano muscular; é uma convenção que combina duração e intensidade relative. A
comparação entre atletas exige cautela, porque o FTP e a resposta individual ao
mesmo estímulo variam. A comparação longitudinal do mesmo atleta, com o mesmo
método, é o uso mais defensável.

\subsection{TRIMP}
\label{sec:trimp}

O \textit{Training Impulse} (TRIMP) de Banister usa a frequência cardíaca para
estimar a carga cardiovascular. A formulação adotada é

\begin{equation}\label{eq:trimp}
TRIMP = t_{\min} \cdot HRR \cdot 0{,}64 \cdot e^{1{,}92\,HRR},
\qquad
HRR = \frac{HR_{\text{média}} - HR_{\text{repouso}}}
{HR_{\text{máx}} - HR_{\text{repouso}}},
\end{equation}

\noindent em que $t_{\min}$ é a duração em minutos e $HRR$ é a fração da
reserva de frequência cardíaca. Os coeficientes 0,64 e 1,92 correspondem à
versão masculina da fórmula de Banister \cite{banister1991,morton1990}. O
\argus{} usa esses coeficientes como base genérica, sem seleção de sexo,
individualização ou calibração. A média da frequência cardíaca é calculada
apenas sobre amostras maiores que zero; quando não há frequência de repouso
válida, o código adota 60 bpm como valor de segurança.

O TRIMP é útil para comparar sessões do mesmo atleta quando a frequência
cardíaca é o sinal disponível. Ele não captura a distribuição de intensidade
dentro da sessão, não distingue intervalos de esforço contínuo com a mesma
média e depende de uma relação linear entre frequência cardíaca e intensidade
que é apenas aproximada. Além disso, os coeficientes foram derivados em uma
população específica; aplicá-los a outro grupo é uma extrapolação.

\subsection{PMC: CTL, ATL e TSB}
\label{sec:pmc}

O \textit{Performance Management Chart} (PMC) transforma o TSS diário em três
curvas: carga crônica (CTL), carga aguda (ATL) e balanço (TSB). O \argus{} usa
médias exponenciais com constantes de 42 e 7 dias:

\begin{equation}\label{eq:ctl}
CTL_n = CTL_{n-1} + \frac{TSS_n - CTL_{n-1}}{42},
\end{equation}

\begin{equation}\label{eq:atl}
ATL_n = ATL_{n-1} + \frac{TSS_n - ATL_{n-1}}{7},
\end{equation}

\begin{equation}\label{eq:tsb}
TSB = CTL - ATL.
\end{equation}

A implementação agrupa as atividades por dia, percorre o calendário do primeiro
registro até a data atual e calcula as médias. O TSB exibido é a diferença entre
CTL e ATL antes da atualização do dia, o que corresponde, na prática, ao estado
do dia anterior. A interpretação ``fitness'', ``fatigue'' e ``form'' é uma
metáfora útil, não uma medida direta de adaptação. O PMC depende da qualidade
do TSS, da consistência dos dados e da premissa de que a carga pode ser somada
entre sessões de tipos diferentes. É uma ferramenta de visualização de
tendência, não um preditor clínico.

\subsection{Desacoplamento cardiovascular e recuperação da frequência cardíaca}
\label{sec:decoupling-hrr}

O desacoplamento (Pw:HR) compara a relação entre potência e frequência cardíaca
em duas metades de uma sessão. O código descarta os primeiros 10 minutos,
calcula a eficiência $EF = P/HR$ em cada metade e aplica

\begin{equation}\label{eq:decoupling}
Decoupling = \frac{EF_1 - EF_2}{EF_1} \times 100\%.
\end{equation}

\noindent Um valor positivo indica que a frequência cardíaca subiu em relação à
potência na segunda metade. O cálculo exige pelo menos 3600 amostras e ignora
amostras com potência ou frequência iguais a zero. O limiar de 5\% é uma
convenção de treino frequentemente usada para sinalizar deriva
cardiovascular; não identifica causa. Desidratação, calor, fadiga, altitude,
cafeína, medicamentos e erro de medição podem produzir o mesmo efeito. Além
disso, o resultado só faz sentido em sessões de intensidade relativamente
constante; em intervalados, a comparação entre metades é pouco informativa.

A recuperação da frequência cardíaca (HRR) mede a queda após o esforço. A
literatura usa geralmente o pico de frequência cardíaca no fim do exercício e
valores em 1 e 2 minutos \cite{cole1999,lamberts2009}. O \argus{} faz uma
aproximação conservadora: ao iniciar o resfriamento, registra a última
frequência cardíaca da sessão como referência e captura a frequência
instantânea em 60 e 120 segundos. Assim,

\begin{equation}\label{eq:hrr}
HRR_1 = HR_{\text{final}} - HR_{60s},
\qquad
HRR_2 = HR_{\text{final}} - HR_{120s}.
\end{equation}

\noindent Se o pico ocorreu antes do fim da sessão, a referência é menor do que
o pico absoluto e o HRR é subestimado. O campo \code{MaxHR} da atividade
registra o máximo real da sessão, mas não é usado no cálculo do HRR. O
protocolo de resfriamento também não é padronizado: a intensidade e a posição
do corpo durante a recuperação afetam a queda. O HRR é, portanto, um indicador
de tendência, não um teste de aptidão.

\subsection{O que é medido, estimado e inferido}
\label{sec:medido-estimado}

A Tabela~\ref{tab:medido-estimado} sintetiza a natureza de cada número
apresentado pelo sistema. Ela é a referência mais importante deste capítulo
para a leitura dos resultados.

\begin{table}[htbp]
    \centering
    \caption{Natureza das grandezas apresentadas pelo Argus}
    \label{tab:medido-estimado}
    \footnotesize
    \renewcommand{\arraystretch}{1.12}
    \begin{tabular}{@{}>{\raggedright\arraybackslash}p{2.8cm}
        >{\raggedright\arraybackslash}p{2.2cm}
        >{\raggedright\arraybackslash}p{4.2cm}
        >{\raggedright\arraybackslash}p{5.5cm}@{}}
        \toprule
        \textbf{Grandeza} & \textbf{Natureza} & \textbf{Origem} &
        \textbf{Limite principal} \\
        \midrule
        Potência, cadência e FC & Medida & Sensores BLE &
        Precisão do sensor, perda de pacotes e taxa de notificação \\
        Distância e velocidade & Estimada & Modelo físico e potência &
        Parâmetros fixos; não é medida de GPS \\
        Elevação e inclinação & Derivada & GPX e suavização &
        Ruído do arquivo, janela e limite de $\pm 25\%$ \\
        NP, IF e TSS & Calculada & Potência e FTP &
        FTP estimada; amostragem e janela de 30 s \\
        TRIMP & Calculada & FC, repouso e máxima &
        Coeficiente populacional; média da FC \\
        Zonas de potência e FC & Classificada & FTP e FC máxima &
        Fronteiras arbitrárias; LTHR não usado nas zonas de FC \\
        Desacoplamento & Calculada & Potência e FC &
        Só em regime relativamente constante; 5\% é heurístico \\
        HRR & Calculada & FC no fim e em 1--2 min &
        Referência é a última FC, não o pico absoluto \\
        VT1 e VT2 & Anotação & Linhas móveis no gráfico &
        Não há detecção automática \\
        VO\textsubscript{2}max & Não medido & --- &
        Exigiria calorimetria indireta ou protocolo validado \\
        Calorias & Estimada & Potência mecânica &
        Conversão aproximada; não é calorimetria \\
        \bottomrule
    \end{tabular}
    \vspace{0.2cm}
    \par\small Fonte: autoria própria.
\end{table}

\section{Engenharia de software}
\label{sec:engenharia}

\subsection{Clean Architecture e a regra das dependências}
\label{sec:clean-architecture}

A \textit{Clean Architecture}, proposta por Robert C. Martin, organiza o
sistema em camadas concêntricas e estabelece que as dependências de código
apontem para dentro: o núcleo de negócio não conhece banco de dados, interface,
framework ou biblioteca de hardware; os detalhes dependem do núcleo
\cite{martin2017clean}. A regra não é uma prescrição de pastas, mas uma
restrição sobre quem pode importar quem. Ela permite substituir a interface, o
banco ou o transporte BLE sem reescrever as regras de negócio, e permite testar
o núcleo sem infraestrutura.

No \argus{}, a separação é pragmaticamente materializada em pacotes Go. O
pacote de domínio contém entidades e interfaces; o pacote de casos de uso
contém a fachada de armazenamento; os pacotes de serviço e de repositório
contêm adaptadores. O \textit{frontend} acessa o núcleo pelos
\textit{bindings} Wails, sem importar pacotes Go. A
Tabela~\ref{tab:camadas} resume a correspondência.

\begin{table}[htbp]
    \centering
    \caption{Camadas, pacotes e responsabilidades}
    \label{tab:camadas}
    \small
    \begin{tabular}{@{}>{\raggedright\arraybackslash}p{2.6cm}
        >{\raggedright\arraybackslash}p{4.6cm}
        >{\raggedright\arraybackslash}p{6.8cm}@{}}
        \toprule
        \textbf{Camada} & \textbf{Pacotes} & \textbf{Responsabilidade} \\
        \midrule
        Apresentação & \code{frontend/src} & Interface, HUD, mapa, gráficos
        e barramento de telemetria \\
        Aplicação & \code{app.go}, \code{internal/usecase} & Orquestração da
        sessão, fachada de armazenamento e exposição ao \textit{frontend} \\
        Domínio & \code{internal/domain} & Entidades, contratos de
        repositório e porta do rolo \\
        Adaptadores & \code{internal/service}, \code{internal/repository} &
        BLE, simulação, GPX, FIT, IA, Strava, Intervals e SQLite \\
        Infraestrutura & SQLite, arquivos, BLE, rede & Sistemas externos
        acionados pelos adaptadores \\
        \bottomrule
    \end{tabular}
    \vspace{0.2cm}
    \par\small Fonte: autoria própria.
\end{table}

A aderência não é dogmática. A fachada \code{StorageFacade} instancia
repositórios SQLite concretos, e as entidades de domínio carregam tags de
serialização e de ORM. Em termos de importação, o domínio não depende de GORM;
em termos de modelo, ele conhece detalhes de persistência. O \code{app.go}
também concentra muita orquestração, o que reduz a pureza dos casos de uso. O
trabalho reconhece essas concessões: a arquitetura é orientada pela regra das
dependências, não é uma implementação de manual.

\subsection{Padrões de projeto}
\label{sec:padroes}

\textbf{Repository.} O padrão isola o acesso a dados atrás de uma interface
definida no domínio \cite{gof1995}. O \argus{} possui seis repositórios de
dados (perfil, atividades, potência, componentes, IA e eventos) e um
gerenciador de conexões. As implementações SQLite ficam em
\code{internal/repository/sqlite}. A interface permite substituir a
persistência nos testes e manter o código de aplicação alheio ao SQL.

\textbf{Facade.} O \code{StorageFacade} oferece uma API única para dezenas de
operações de armazenamento, incluindo a troca do banco ativo quando o perfil
muda. Sem a fachada, cada controlador precisaria conhecer a ordem de abertura
dos bancos e as regras de migração.

\textbf{Observer.} O \code{TelemetryEventBus}, no \textit{frontend}, publica
cada amostra para HUD, gráficos e mapa sem que o produtor conheça os
consumidores. O padrão desacopla a chegada do dado da renderização; não
significa que a interface atualize a 60 quadros por segundo, nem que os dados
cheguem nessa taxa. A fluidez visual é obtida por atualizações frequentes e,
no mapa, por limitação de taxa e animação de câmera.

\textbf{Strategy/Adapter.} A porta \code{TrainerService} permite trocar a
implementação BLE real por um rolo virtual ou por um \textit{mock} de
teste. O mesmo contrato expõe modos ERG e simulação, de modo que o restante do
sistema não precisa saber se o transporte é FTMS ou FE-C.

\subsection{Local-first e suas fronteiras}
\label{sec:local-first}

O conceito de \textit{local-first} descreve software em que os dados primários
e a capacidade de trabalhar sobre eles permanecem no dispositivo do usuário,
mesmo quando há serviços remotos \cite{kleppmann2019}. O \argus{} adota essa
orientação em um sentido restrito: perfis, atividades, recordes e configurações
ficam em SQLite local; o cálculo fisiológico e a simulação rodam no
dispositivo; as integrações com Strava, Intervals.icu, tiles de mapa e Ollama
são opcionais ou externas. Não há sincronização entre dispositivos, resolução
de conflitos, \textit{CRDT} ou conta obrigatória. O sistema é local-first no
controle do dado, não no sentido completo de colaboração multi-dispositivo.

\section{Inteligência artificial generativa local}
\label{sec:ia-local}

Um modelo de linguagem de grande escala (LLM) pode ser executado localmente por
meio de servidores de inferência como o Ollama \cite{ollama2024}. O \argus{}
usa esse caminho para um assistente de treino: o modelo recebe o perfil do
atleta (FTP, peso, nível, histórico recente e métricas agregadas), interpreta
o pedido em linguagem natural e pode devolver um plano estruturado em JSON, que
o sistema converte em ZWO. O endereço padrão é \texttt{127.0.0.1:11434} e o
modelo padrão é \code{qwen2.5:3b}; o usuário pode trocar o modelo.

A vantagem é a privacidade: os dados biométricos não precisam sair do
dispositivo. O limite é a qualidade do modelo e a ausência de validação
fisiológica. Um LLM pequeno pode alucinar durações, intensidades ou relações
entre FTP e zonas; pode produzir JSON sintaticamente válido e fisiologicamente
inadequado. O \argus{} valida a estrutura do JSON e converte os campos
conhecidos, mas não verifica se o plano é seguro, progressivo ou adequado ao
atleta. O assistente é uma ferramenta de ideação e edição, não um treinador
certificado. A recomendação de revisão humana permanece.

\section{Síntese do capítulo}
\label{sec:sintese-fundamentacao}

A Tabela~\ref{tab:conceitos-sintese} fecha o capítulo relacionando os conceitos
centrais às decisões do \argus{} e aos seus limites. A leitura conjunta das
tabelas~\ref{tab:medido-estimado} e~\ref{tab:conceitos-sintese} evita duas
leituras equivocadas: a de que o sistema mede tudo o que exibe, e a de que a
simulação é uma cópia da estrada.

\begin{table}[htbp]
    \centering
    \caption{Síntese conceitual do capítulo}
    \label{tab:conceitos-sintese}
    \footnotesize
    \renewcommand{\arraystretch}{1.1}
    \begin{tabular}{@{}>{\raggedright\arraybackslash}p{3.2cm}
        >{\raggedright\arraybackslash}p{5.2cm}
        >{\raggedright\arraybackslash}p{5.2cm}@{}}
        \toprule
        \textbf{Conceito} & \textbf{Decisão no Argus} & \textbf{Limite assumido} \\
        \midrule
        BLE/GATT & Assinar HR, potência e FTMS; timestamp na chegada &
        Taxa não é fixa; não há reamostragem \\
        FTMS e FE-C & Dois caminhos de controle no desktop; FE-C no Android &
        FE-C sobre BLE não é ANT+ nativo \\
        GPX & Perfil, inclinação e distância 3D no desktop &
        Relevo aproximado; sem superfície, clima ou curvas \\
        FIT & Exportação binária interoperável &
        Aceitação não prova exatidão; validação formal pendente \\
        Martin & Modelo quase estático com parâmetros fixos &
        Sem inércia, vento, frenagem ou dinâmica do rolo \\
        Fisiologia & NP, IF, TSS, TRIMP, PMC, HRR e desacoplamento &
        Estimativas e convenções; sem validação clínica \\
        Clean Architecture & Domínio com interfaces e adaptadores &
        Concessões pragmáticas em \code{app.go} e \code{StorageFacade} \\
        IA local & Ollama, prompt com perfil e geração de ZWO &
        Modelo pode alucinar; não há validação fisiológica \\
        \bottomrule
    \end{tabular}
    \vspace{0.2cm}
    \par\small Fonte: autoria própria.
\end{table}

O próximo capítulo descreve como essas decisões foram organizadas em uma
arquitetura executável e quais compensações foram feitas para sustentar o
\textit{desktop} e o Android no mesmo projeto.
